\documentclass[12pt]{article}
\usepackage[utf8]{inputenc}
\usepackage[T1]{fontenc}
\usepackage{hyperref}
\title{Paiwan Language Resources for K-12 Curriculum Development}
\author{Compiled by Claude Sonnet 5.5 (claude-sonnet-5-5) for Andr\'as Kornai}
\date{}
\begin{document}
\maketitle

\section{Introduction}
Paiwan (ISO 639-3 pwn) is a Formosan language of the Austronesian family spoken in southern Taiwan. Sources found in search disagree on speaker numbers (roughly 66,000 to 96,000), so the figure is uncertain. Since 2017 Taiwan has treated indigenous languages as national languages and included them in the school curriculum, and Paiwan is taught in some schools in northern and southern varieties. Web search surfaced few directly downloadable Paiwan teaching materials; this dossier is correspondingly short.

\section{Primary Educational Resources}
\begin{itemize}
\item New Taipei City course outline for Paiwan (``300 words and their derivatives'' from the official indigenous-language learning word list): \url{https://www.buda.ntpc.gov.tw/plugin/course_information/uploadfiles/20240403115048331183207.pdf}
\item Taipei primary-school curriculum plans that distinguish Northern and Southern Paiwan: \url{https://www.cp.ptc.edu.tw/storage/134787/134787_114_B-02_2A.pdf}
\end{itemize}
Search results did not reveal the titles or links of the national graded Paiwan textbooks themselves; these should be sought from the Ministry of Education and the National Academy for Educational Research.

\section{Linguistic and Reference Materials}
\begin{itemize}
\item Wikipedia, Paiwan language (overview, dialects): \url{https://en.wikipedia.org/wiki/Paiwan_language}
\item Wikipedia, Formosan languages: \url{https://en.wikipedia.org/wiki/Formosan_languages}
\item Handbook of Formosan Languages (Brill): \url{https://referenceworks.brillonline.com/browse/handbook-of-formosan-languages-the-indigenous-languages-of-taiwan/alpha/p}
\item A comparative study on Formosan phonology: Paiwan and Budai Rukai (UT Austin repository): \url{https://repositories.lib.utexas.edu/items/aaf7ccba-4293-4498-ba57-ac2ee8357e30/full}
\item The Language Closet, Languages of Taiwan: Paiwan: \url{https://thelanguagecloset.com/2020/12/19/languages-of-taiwan-paiwan-vinuculjan-pinayuanan/}
\item Rosetta Project entry for pwn: \url{https://rosettapanglossia.longnow.org/index.php?title=Pwn}
\end{itemize}

\section{Religious and Community Resources}
No Paiwan religious or community-organisation resources were identified in the search results.

\section{Digital and Audio Resources}
\begin{itemize}
\item The Council of Indigenous Peoples' online dictionary of indigenous languages is mentioned in search results as covering 14 peoples and 16 languages; I did not see a direct URL or confirm Paiwan coverage, so it is not listed as a resource.
\end{itemize}

\section{Curriculum Development Context}
Research on Eastern Paiwan textbooks (gender portrayal) was presented at an IAFOR conference: \url{https://egen.iafor.org/presentation/submission110449/}. Paiwan has two main teaching varieties (Northern and Southern) so a curriculum must choose or accommodate both. The existing school provision is mostly Mandarin-medium with Paiwan as a subject, and little beyond that was found.

\end{document}